\documentclass[11pt,a4paper]{report}

\usepackage[T1]{fontenc}
\usepackage[utf8]{inputenc}
\usepackage[french]{babel}
\usepackage{lmodern}
\usepackage[a4paper,margin=2.2cm]{geometry}
\usepackage{xcolor}
\usepackage{titlesec}
\usepackage{tabularx}
\usepackage{booktabs}
\usepackage{array}
\usepackage{enumitem}
\usepackage{fancyhdr}
\usepackage{longtable}
\usepackage{needspace}
\usepackage[hidelinks]{hyperref}

% ---- Couleurs de marque ----
\definecolor{sagiteal}{HTML}{00745A}
\definecolor{sagidark}{HTML}{1A1A2E}
\definecolor{sagigray}{HTML}{6B7280}
\definecolor{sagibg}{HTML}{F4F6F8}
\definecolor{sagired}{HTML}{B91C1C}
\definecolor{sagiamber}{HTML}{B45309}

\hypersetup{
  pdftitle={SAGI SCHOOL -- Reprise des donnees a la migration},
  pdfauthor={HADY GESMAN},
  colorlinks=false,
}

% ---- Titres ----
\titleformat{\chapter}[hang]
  {\bfseries\Large\color{sagiteal}}{\thechapter.}{0.6em}{}
\titlespacing{\chapter}{0pt}{0pt}{12pt}
\titleformat{\section}
  {\bfseries\large\color{sagidark}}{\thesection}{0.6em}{}
\titleformat{\subsection}
  {\bfseries\normalsize\color{sagidark}}{\thesubsection}{0.6em}{}

% ---- En-tête / pied ----
\pagestyle{fancy}
\fancyhf{}
\renewcommand{\headrulewidth}{0.4pt}
\renewcommand{\footrulewidth}{0.4pt}
\fancyhead[L]{\small\color{sagiteal}\textbf{SAGI SCHOOL}}
\fancyhead[R]{\small\color{sagigray}Reprise des données à la migration}
\fancyfoot[L]{\small\color{sagigray}by HADY GESMAN}
\fancyfoot[C]{\small\color{sagigray}Conforme SYSCOHADA Révisé}
\fancyfoot[R]{\small\thepage}

\setlist[itemize]{leftmargin=1.3em,itemsep=2pt,topsep=3pt}
\setlist[enumerate]{leftmargin=1.5em,itemsep=3pt,topsep=3pt}
\renewcommand{\tabularxcolumn}[1]{>{\small}m{#1}}

% ---- Encadrés ----
\newcommand{\encadre}[2]{%
  \begingroup
  \setlength{\fboxsep}{9pt}%
  \noindent\colorbox{sagibg}{%
    \begin{minipage}{\dimexpr\linewidth-2\fboxsep\relax}
    \textbf{\color{#1}#2}
    \end{minipage}}%
  \par\vspace{8pt}
  \endgroup}

\newcommand{\pointcle}[1]{\encadre{sagiteal}{#1}}
\newcommand{\alerte}[1]{\encadre{sagired}{#1}}

\begin{document}

% =========================== PAGE DE TITRE ===========================
\begin{titlepage}
\centering
\vspace*{2.5cm}
{\Huge\bfseries\color{sagiteal} SAGI SCHOOL\par}
\vspace{0.4cm}
{\Large\color{sagidark} Système de Gestion Scolaire\par}
\vspace{0.2cm}
{\large\color{sagigray} \textit{by HADY GESMAN}\par}
\vspace{2.2cm}
{\LARGE\bfseries\color{sagidark} Reprise des données à la migration\par}
\vspace{0.5cm}
{\large\color{sagigray} Rapport de développement --- 4 lots livrés\par}
\vspace{0.3cm}
{\normalsize\color{sagigray} Cas de référence : Complexe Shoumoul Excellence\par}
\vspace{2.6cm}

\begin{tabular}{rl}
\textbf{Date} & 27 juillet 2026\\[4pt]
\textbf{Branche} & \texttt{feat/impaye-anterieur-migration}\\[4pt]
\textbf{Commits} & 4\\[4pt]
\textbf{Portée} & 31 fichiers, +3 276 / --57 lignes\\[4pt]
\textbf{Tests} & 123 verts (82 nouveaux)\\[4pt]
\textbf{Statut} & \textbf{\color{sagiamber}En attente de validation}\\[4pt]
\end{tabular}

\vfill
{\small\color{sagigray} Document interne --- pour décision de mise en production\par}
\end{titlepage}

% Sommaire tenu sur une page : chapitres et sections seulement, corps réduit.
{\small
\setcounter{tocdepth}{1}
\tableofcontents
}
\newpage

% =========================== RÉSUMÉ EXÉCUTIF ===========================
\chapter*{Résumé exécutif}
\addcontentsline{toc}{chapter}{Résumé exécutif}

\section*{Le problème}

L'installation de SAGI SCHOOL au \textbf{Complexe Shoumoul Excellence} était
bloquée. Le logiciel savait calculer ce que chaque élève doit pour l'année en
cours, à partir des tarifs paramétrés par niveau. Il ne savait pas enregistrer
ce que les familles devaient \emph{déjà} en arrivant.

Or c'est la situation normale sur le terrain : une école qui bascule sur un
logiciel traîne des années d'ardoises, connues de manière approximative, sans
détail exploitable poste par poste. Le directeur sait dire « cette famille doit
45 000 francs », rarement « dont 12 000 d'inscription 2023 et 33 000 de
mensualités ». Exiger le détail, c'est bloquer la migration.

Deux besoins connexes sont apparus au même moment : constituer une base
d'élèves fiable année après année, et retrouver le dossier complet d'un ancien
élève longtemps après son départ.

\section*{Ce qui a été livré}

Quatre lots, développés et testés :

\begin{enumerate}
\item \textbf{Impayé antérieur saisissable} --- un montant global par élève
      suffit, saisissable de trois façons, avec l'écriture comptable
      correspondante. \emph{Lot qui débloque Shoumoul.}
\item \textbf{Matricule de promotion} --- le matricule dit désormais l'année
      et le rang d'entrée de l'élève, et ne change plus jamais.
\item \textbf{Parcours scolaire et anciens élèves} --- l'histoire complète d'un
      enfant, de son entrée à sa sortie, et une base historique des diplômés.
\item \textbf{Santé de la migration} --- un tableau de bord de ce qui reste à
      compléter dans les données reprises.
\end{enumerate}

\section*{Principe directeur}

\pointcle{La migration n'est jamais bloquante. Une école peut installer le
logiciel sans aucune donnée financière, travailler dès le premier jour, et
compléter progressivement --- élève par élève, sans jamais refaire l'import.}

\section*{Décision demandée}

Autorisation de pousser les 4 commits, d'ouvrir la demande de fusion, de
publier une version et de déployer. Le plan de mise en production détaillé
figure au chapitre~\ref{chap:deploiement}.

\newpage

% =========================== CHAPITRE 1 ===========================
\chapter{Le contexte : pourquoi la migration bloquait}

\section{Ce que le logiciel savait déjà faire}

Avant ces travaux, SAGI SCHOOL gérait correctement le cycle normal d'une
école :

\begin{itemize}
\item tarifs paramétrés par niveau (inscription, mensualité, uniforme,
      fournitures, services optionnels) ;
\item calcul automatique du total dû par élève, au prorata de sa date
      d'entrée et de sa prise en charge sociale éventuelle ;
\item suivi des paiements, alertes de retard, reçus, comptabilité SYSCOHADA ;
\item report automatique des impayés d'un exercice sur le suivant, lors de la
      clôture.
\end{itemize}

Ce dernier point mérite attention : le mécanisme de report existait, avec toute
sa mécanique comptable. Mais il supposait que \textbf{l'année précédente existe
dans le logiciel}. Pour une école qui migre, c'est précisément ce qui manque.

\section{Le blocage constaté chez Shoumoul}

Les élèves du Complexe Shoumoul Excellence arrivaient avec des impayés nés
avant l'installation. Le logiciel n'avait aucun endroit où les enregistrer :

\begin{itemize}
\item le total dû affiché était \emph{faux}, car amputé de la dette ancienne ;
\item l'école ne pouvait pas relancer les familles sur le bon montant ;
\item la comptabilité ne portait pas ces créances au bilan.
\end{itemize}

\section{La difficulté de fond}

\alerte{Les données financières fiables sont rares sur le terrain. Un système
qui exige le détail poste par poste d'une dette de trois ans ne sera jamais
rempli --- l'école abandonnera, ou saisira n'importe quoi.}

C'est cette contrainte qui a dicté la conception : accepter le niveau de
précision que l'école peut réellement fournir, et permettre d'affiner ensuite.

\section{Les quatre niveaux de précision acceptés}

\begin{center}
\begin{tabularx}{\linewidth}{@{}p{0.5cm}p{4.2cm}X@{}}
\toprule
\textbf{N\textsuperscript{o}} & \textbf{Niveau} & \textbf{Ce que l'école doit savoir dire} \\
\midrule
1 & Rien & Aucune donnée financière. L'élève démarre à zéro payé. \\
2 & Dette actuelle & « Cette famille doit 30 000 francs à ce jour. » \\
3 & Détail par poste & Inscription réglée, 3 mensualités payées, etc. \\
4 & + Impayé antérieur & « Et elle devait déjà 45 000 avant. » \emph{(nouveau)} \\
\bottomrule
\end{tabularx}
\end{center}

Les trois premiers niveaux existaient. Le quatrième est ajouté, et il est
\textbf{indépendant} : il se cumule avec n'importe lequel des autres, ou se
remplit seul.

% =========================== CHAPITRE 2 ===========================
\chapter{Lot 1 --- Impayé antérieur saisissable}

\section{Le principe retenu}

Un seul montant global par élève, plus une note libre sur son origine
(« 2024--2025 », « ardoise cahier », « ancien logiciel »). Aucune ventilation
par poste n'est demandée --- c'est précisément ce que les écoles ne peuvent
pas fournir.

\section{Trois points d'entrée}

\begin{center}
\begin{tabularx}{\linewidth}{@{}p{4.2cm}X@{}}
\toprule
\textbf{Point d'entrée} & \textbf{Usage} \\
\midrule
Fiche élève & Cas par cas, à tout moment de l'année. \\
Import Excel & Colonnes « Impayé antérieur » et « Origine », avec
  tolérance sur les intitulés (\textit{ardoise}, \textit{arriérés},
  \textit{dette antérieure}, \textit{reliquat}\ldots). \\
Saisie en lot & Une grille avec recherche et total en direct, pour traiter
  plusieurs centaines d'élèves en une passe. \\
\bottomrule
\end{tabularx}
\end{center}

\section{Traitement comptable}

L'impayé antérieur est une créance née \emph{avant} l'exercice en cours. Le
produit correspondant a été constaté à l'époque, dans l'ancien système ou sur
les cahiers de l'école.

\pointcle{Écriture passée : 411 Débit / 890 Crédit, à la date d'ouverture de
l'exercice. Aucun compte de produit (706), aucun mouvement de trésorerie.}

Cette règle est essentielle : rejouer le produit gonflerait le résultat de
l'année en cours avec des frais qui ne lui appartiennent pas, et fausserait le
compte de résultat comme le bilan.

\section{Deux choix de conception qui comptent}

\subsection{Une écriture auto-réparatrice}

Une migration se nettoie progressivement : un montant sera corrigé plusieurs
fois. Plutôt que d'empiler des écritures d'ajustement, le système
\textbf{recalcule l'écriture en entier} à chaque modification.

Conséquences concrètes :

\begin{itemize}
\item trois corrections successives ne laissent \textbf{qu'une seule pièce} au
      journal, celle du dernier montant ;
\item le numéro de pièce reste stable d'une correction à l'autre ;
\item remettre le montant à zéro fait disparaître l'écriture, sans trace
      parasite.
\end{itemize}

Cette approche découle directement d'un incident antérieur : un mécanisme
d'ajustement empilé avait fini par ronger les produits d'une école migrée
jusqu'à afficher zéro recette, sans qu'aucun contrôle ne le signale.

\subsection{Un refus ligne par ligne, jamais en bloc}

La saisie en lot applique ce qu'elle peut et rejette le reste avec le motif
exact. Sur trois cents élèves saisis à la main, tout annuler pour une faute de
frappe serait le meilleur moyen de décourager l'école.

\section{Garde-fous}

\begin{itemize}
\item montant négatif refusé ;
\item exercice clôturé : modification refusée (une année fermée reste en
      lecture seule) ;
\item montant inférieur à ce qui a déjà été encaissé sur cette dette : refusé,
      avec le montant déjà réglé dans le message.
\end{itemize}

\section{Affichage}

La colonne principale de la liste des élèves affiche désormais le \textbf{dû
global} --- année en cours plus ardoise ancienne --- avec la décomposition au
survol et une colonne « dont antérieur » à côté.

Le niveau d'alerte (Attention / Urgent / Critique), lui, continue de ne juger
que l'année en cours. Sans cette précaution, toute l'école basculerait en
« Critique » le jour de la migration, rendant l'indicateur inutilisable.

% =========================== CHAPITRE 3 ===========================
\chapter{Lot 2 --- Matricule de promotion}

\section{Le défaut de l'ancien format}

L'ancien matricule avait la forme \texttt{AAAA-CODE-NNNNNN}, où :

\begin{itemize}
\item \texttt{AAAA} était l'\textbf{année civile du jour de la saisie} ;
\item \texttt{NNNNNN} un compteur global à l'école.
\end{itemize}

Deux élèves d'une même promotion, l'un inscrit en septembre et l'autre en
janvier, portaient donc deux années différentes. Et le numéro ne disait rien de
la cohorte. Impossible d'en tirer une base ordonnée des premiers aux derniers
inscrits.

\section{Le nouveau format}

\pointcle{\texttt{2025-SHE-0042} se lit : quarante-deuxième élève entré en
2025--2026 au Complexe Shoumoul Excellence.}

\begin{itemize}
\item \texttt{2025} : année de début de l'exercice d'entrée --- la promotion ;
\item \texttt{SHE} : code de l'établissement ;
\item \texttt{0042} : rang dans la promotion, par ordre d'entrée.
\end{itemize}

Un élève inscrit en janvier reste de la promotion de son exercice. Le matricule
est attribué une fois et suit l'enfant jusqu'à sa sortie : réinscription,
clôture d'exercice, changement de niveau ne le modifient jamais.

\section{Deux dates, deux usages}

Un point technique aux conséquences réelles a été découvert en explorant ce
lot.

\alerte{La date d'entrée réelle des élèves était déjà perdue. Le mécanisme de
réinscription repositionne chaque année la date d'inscription au début du
nouvel exercice, pour le calcul du prorata des mensualités. Dès la deuxième
année, plus aucune trace de la date d'arrivée effective.}

Deux champs figés à vie ont donc été ajoutés, recopiés à chaque
réinscription :

\begin{center}
\begin{tabularx}{\linewidth}{@{}p{4.5cm}X@{}}
\toprule
\textbf{Champ} & \textbf{Rôle} \\
\midrule
Date d'inscription & Repositionnée chaque année. Sert au prorata des
  mensualités dues. \\
Date d'entrée \emph{(nouveau)} & Ne bouge jamais. Première arrivée dans
  l'établissement. \\
Année d'entrée \emph{(nouveau)} & La promotion. \\
\bottomrule
\end{tabularx}
\end{center}

\section{Reprise de l'existant}

Une commande de renumérotation traite les bases déjà installées. Elle
fonctionne en \textbf{diagnostic par défaut} : elle affiche la correspondance
ancien~$\rightarrow$~nouveau sans rien écrire. L'écriture demande une option
explicite.

Trois précautions ont été nécessaires :

\begin{enumerate}
\item \textbf{Regrouper les fiches par enfant réel.} Un même enfant possède une
      fiche par année. Toutes doivent recevoir le même matricule, sinon on
      casse justement le suivi qu'on cherche à établir. Le regroupement combine
      le lien de réinscription \emph{et} l'égalité de matricule --- une fiche
      importée avant le report n'est pas liée, et une base ancienne peut
      manquer de matricule.
\item \textbf{Écrire en deux passes.} Quand deux élèves échangent leurs rangs,
      écrire directement violerait la contrainte d'unicité en cours de route.
      Les matricules sont donc d'abord parqués sur une valeur temporaire.
\item \textbf{Conserver l'ancien matricule.} Il est stocké à part, et la
      recherche accepte les deux numéros --- les carnets papier et les anciens
      reçus de l'école restent exploitables.
\end{enumerate}

La commande est rejouable : un second passage recalcule les mêmes matricules et
ne touche plus rien.

\alerte{Point d'attention pour le déploiement : tant que cette commande n'a pas
été lancée, l'année et la date d'entrée restent vides sur les fiches
existantes. C'est un choix assumé --- un remplissage automatique prendrait la
promotion de la fiche courante au lieu de celle d'entrée pour les élèves
réinscrits, donnant une valeur \emph{fausse} plutôt qu'une valeur vide sur une
base censée devenir la référence historique.}

% =========================== CHAPITRE 4 ===========================
\chapter{Lot 3 --- Parcours scolaire et anciens élèves}

\section{Le constat}

Un enfant possède une fiche par exercice. Son histoire était donc éparpillée
sur autant de lignes qu'il avait passé d'années dans l'établissement, sans
aucun écran ni aucune interface pour la rassembler.

\section{Le parcours}

Un onglet « Parcours scolaire », ouvrable depuis la fiche élève ou depuis la
base des anciens, présente la scolarité complète année par année : classe,
statut, dû, payé, reste, reliquat repris. Il est complété par une synthèse
depuis l'entrée et par un \textbf{dossier de scolarité en PDF}, à remettre à
une famille qui part ou à ressortir des archives.

\subsection{Un piège de calcul évité}

\alerte{La dette totale d'un parcours n'est pas la somme des restes annuels.}

Un impayé est reconduit en reliquat chaque année. L'additionner sur toutes les
lignes reviendrait à compter la même ardoise autant de fois qu'elle a traversé
d'exercices : sur trois ans, 50 000 francs deviendraient 150 000. C'est le dû
de la dernière année qui fait foi, car il porte déjà tout l'historique. Un test
automatisé verrouille ce comportement.

\section{La base des anciens élèves}

Un écran dédié, \textbf{indépendant de l'exercice affiché} : on doit y
retrouver un diplômé de 2019 comme un transféré de l'an dernier.

\begin{itemize}
\item une ligne par enfant, pas par fiche ;
\item diplômés, transférés et abandons, filtrables ;
\item recherche par nom, matricule ou ancien matricule ;
\item pour chacun : promotion, date d'entrée, dernière année, nombre d'années,
      total réglé sur toute la scolarité, et solde éventuellement encore dû ;
\item total des créances des anciens en indicateur.
\end{itemize}

Le statut retenu est celui de la dernière fiche. Un enfant réinscrit après un
abandon n'est donc plus un ancien --- il est redevenu élève.

\section{Le sort des sortants endettés}

Un diplômé qui part en devant de l'argent posait un dilemme.

\begin{itemize}
\item \textbf{Le réinscrire} comme élève de l'année suivante, comme le faisait
      le système : il pollue la liste et les listes d'appel.
\item \textbf{Ne rien créer du tout} : sa créance disparaît du bilan du nouvel
      exercice --- inacceptable comptablement.
\end{itemize}

La solution retenue crée une \textbf{fiche de créance} : elle porte l'écriture
de report 411/890, donc la créance reste au bilan, mais elle est tenue hors de
la liste des élèves et hors des effectifs. Elle se consulte depuis la base des
anciens élèves.

% =========================== CHAPITRE 5 ===========================
\chapter{Lot 4 --- Santé de la migration}

\section{Pourquoi ce tableau}

Une migration ne se termine pas le jour de l'installation. L'école complète
ensuite --- encore faut-il qu'elle sache ce qui reste à compléter. Sans ce
tableau, les trous se découvrent six mois plus tard, au moment d'éditer un
bilan.

\section{Les six contrôles}

\begin{center}
\begin{tabularx}{\linewidth}{@{}p{4.6cm}X@{}}
\toprule
\textbf{Contrôle} & \textbf{Ce qu'il détecte} \\
\midrule
Situation financière & Élèves sans aucun paiement ni impayé antérieur
  enregistré : ceux dont on ne sait rien. \\
Impayés antérieurs & Nombre d'ardoises saisies et montant total. \\
Identité d'entrée & Matricule, année ou date d'entrée manquants : la commande
  de renumérotation n'a pas été lancée. \\
Sections sans tarif & Une section sans frais renseignés : tous ses élèves
  doivent zéro franc, et rien ne le signale. \\
Élèves sans section & Leur dû ne peut pas être calculé du tout. \\
Équilibre du journal & Écart entre débits et crédits : une reprise ou un
  recalage a mal tourné. \\
\bottomrule
\end{tabularx}
\end{center}

Quatre de ces six contrôles visent des anomalies \textbf{parfaitement
silencieuses} : rien, dans l'usage quotidien, ne les révèle.

\section{Deux partis pris}

\begin{itemize}
\item \textbf{Aucun contrôle n'est bloquant.} Trois niveaux seulement : correct,
      information, à traiter. C'est une liste de courses, pas un barrage.
\item \textbf{Le tableau ne corrige rien.} Il compte, et dit où aller.
\end{itemize}

% =========================== CHAPITRE 6 ===========================
\chapter{Anomalies corrigées en chemin}

Trois défauts ont été trouvés pendant les travaux. Aucun n'était visible dans
l'usage courant.

\section{Messages d'erreur avalés à l'enregistrement d'une fiche}

L'enregistrement d'un élève affichait « impossible d'ajouter » quelle que soit
la cause réelle du refus. Les nouveaux garde-fous auraient donc laissé l'école
corriger à l'aveugle. Le motif exact du serveur est désormais remonté à
l'utilisateur.

\section{Date d'entrée perdue dès la deuxième année}

Détaillé au chapitre~3 : la réinscription écrasait chaque année la date
d'arrivée réelle. Le défaut existait depuis la mise en place du report des
reliquats. Il est corrigé par le champ figé de date d'entrée.

\section{Scolarité facturée à un enfant déjà parti}

Découvert en écrivant les tests du lot 3. Une fiche de créance --- ouverte pour
suivre l'ardoise d'un élève sorti --- réclamait \emph{aussi} la scolarité
complète de l'année, soit 300 000 francs facturés à un enfant qui ne remettra
pas les pieds dans l'établissement. Le total attendu de ces fiches est
désormais nul.

% =========================== CHAPITRE 7 ===========================
\chapter{Qualité et couverture de tests}

\section{Chiffres}

\begin{center}
\begin{tabularx}{\linewidth}{@{}Xr@{}}
\toprule
\textbf{Indicateur} & \textbf{Valeur} \\
\midrule
Commits & 4 \\
Fichiers touchés & 31 (16 créés, 15 modifiés) \\
Lignes & +3 276 / --57 \\
Tests automatisés sur le périmètre & 123, tous verts \\
\quad dont nouveaux & 82 \\
Migrations de base de données & 3 \\
Nouveaux points d'API & 5 \\
Compilation de l'interface & sans erreur \\
\bottomrule
\end{tabularx}
\end{center}

\section{Répartition des nouveaux tests}

\begin{center}
\begin{tabularx}{\linewidth}{@{}Xr@{}}
\toprule
\textbf{Module} & \textbf{Tests} \\
\midrule
Impayé antérieur (écritures, auto-réparation, garde-fous, API) & 22 \\
Matricule de promotion et renumérotation & 21 \\
Parcours, anciens élèves, sortants endettés, PDF & 20 \\
Santé de la migration & 15 \\
Import Excel --- colonne impayé antérieur & 4 \\
\bottomrule
\end{tabularx}
\end{center}

\section{Ce que les tests protègent en priorité}

\begin{itemize}
\item l'absence de produit (706) sur les reprises de dette --- sans quoi le
      résultat de l'exercice serait gonflé ;
\item l'équilibre du journal après chaque opération ;
\item le caractère auto-réparateur des écritures, y compris après plusieurs
      corrections successives et une remise à zéro ;
\item la non-sommation des restes annuels dans un parcours ;
\item l'isolation entre écoles : aucune opération ne doit franchir la frontière
      d'un établissement ;
\item la survie de l'identité d'entrée à la réinscription.
\end{itemize}

% =========================== CHAPITRE 8 ===========================
\chapter{Plan de mise en production}
\label{chap:deploiement}

\section{État actuel}

Les 4 commits sont sur une branche locale. \textbf{Rien n'a été poussé}, aucune
demande de fusion n'est ouverte, aucune version n'est publiée.

\section{Séquence proposée après validation}

\begin{enumerate}
\item \textbf{Pousser} la branche et ouvrir la demande de fusion vers la
      branche principale.
\item \textbf{Fusionner} après relecture.
\item \textbf{Publier une version} par étiquette Git, ce qui déclenche la
      chaîne d'intégration continue et la production de l'installeur Windows.
\item \textbf{Déployer le cloud} : construction de l'interface en local, envoi
      vers le serveur, purge du cache du réseau de diffusion.
\item \textbf{Déployer chez Shoumoul} (installation locale Windows) selon la
      procédure ci-dessous.
\end{enumerate}

\section{Procédure sur l'installation Shoumoul}

\begin{enumerate}
\item \textbf{Sauvegarde de la base} avant toute opération.
\item \textbf{Appliquer les 3 migrations} de base de données.
\item \textbf{Lancer la renumérotation des matricules en mode diagnostic} et
      \emph{examiner la sortie} : la correspondance ancien~$\rightarrow$~nouveau
      doit être cohérente, notamment pour les élèves présents depuis plusieurs
      années.
\item \textbf{Valider avec l'école}, puis relancer la commande avec l'option
      d'écriture.
\item \textbf{Ouvrir l'onglet « Santé de la migration »} : les contrôles
      « identité d'entrée » et « équilibre du journal » doivent être au vert.
\item \textbf{Saisir les impayés antérieurs} avec la direction de l'école, via
      la grille de saisie en lot. Les montants doivent être validés par
      l'établissement avant enregistrement.
\end{enumerate}

\alerte{Aucune de ces opérations ne modifie un exercice clôturé. Les écritures
de reprise sont passées exclusivement dans l'exercice actif, conformément à la
règle appliquée depuis la mise en place du report des reliquats.}

\section{Retour arrière}

Les trois migrations sont additives : elles ajoutent des colonnes sans en
supprimer ni en transformer. Un retour à la version précédente du logiciel
laisse les données intactes.

La renumérotation des matricules, elle, écrit dans la base --- d'où la
sauvegarde préalable et l'obligation de passer par le diagnostic. L'ancien
matricule reste conservé dans tous les cas.

% Garde la section entière sur une même page : sans cela la dernière puce
% se retrouve seule sur une page blanche.
\Needspace*{8\baselineskip}
\section{Points à surveiller après déploiement}

\begin{itemize}
\item la lecture du dû global par l'école : la colonne principale a changé de
      signification, elle inclut désormais l'ardoise ancienne ;
\item les fiches de créance des élèves sortis, qui n'apparaissent plus dans la
      liste des élèves --- comportement voulu, à expliquer à la direction ;
\item le total des créances anciennes, nouvel indicateur de la base des
      anciens élèves.
\end{itemize}

% =========================== ANNEXE ===========================
\appendix
\chapter{Inventaire technique}

\section{Modules créés}

\begin{center}
\begin{tabularx}{\linewidth}{@{}>{\small\ttfamily}p{6.2cm}X@{}}
\toprule
\textbf{Fichier} & \textbf{Rôle} \\
\midrule
paiements/reliquat\_migration.py & Saisie et écriture de l'impayé
  antérieur, synchronisation auto-réparatrice. \\
eleves/matricules.py & Attribution du matricule de promotion. \\
eleves/rebasage.py & Renumérotation de l'existant. \\
eleves/parcours.py & Parcours d'un élève, base des anciens,
  regroupement des fiches par enfant. \\
eleves/sante\_migration.py & Diagnostic de complétude. \\
commands/rebaser\_matricules.py & Commande d'administration. \\
templates/pdf/parcours\_eleve.html & Dossier de scolarité PDF. \\
\bottomrule
\end{tabularx}
\end{center}

\section{Nouveaux points d'API}

\begin{center}
\begin{tabularx}{\linewidth}{@{}>{\small\ttfamily}p{6.8cm}X@{}}
\toprule
\textbf{Route} & \textbf{Usage} \\
\midrule
GET/POST /eleves/impayes-anterieurs/ & Grille de saisie en lot. \\
\texttt{GET /eleves/\{id\}/parcours/} & Scolarité complète. \\
\texttt{GET /eleves/\{id\}/parcours-pdf/} & Dossier de scolarité. \\
GET /eleves/anciens/ & Base historique des sortis. \\
GET /eleves/sante-migration/ & Diagnostic de migration. \\
\bottomrule
\end{tabularx}
\end{center}

\section{Migrations de base de données}

\begin{center}
\begin{tabularx}{\linewidth}{@{}p{2.2cm}X@{}}
\toprule
\textbf{Numéro} & \textbf{Contenu} \\
\midrule
\texttt{0021} & Note d'origine de l'impayé antérieur. \\
\texttt{0022} & Année d'entrée, date d'entrée, ancien matricule. \\
\texttt{0023} & Indicateur de fiche de créance. \\
\bottomrule
\end{tabularx}
\end{center}

\section{Historique des commits}

\begin{center}
\begin{tabularx}{\linewidth}{@{}p{2.2cm}X@{}}
\toprule
\textbf{Référence} & \textbf{Objet} \\
\midrule
\texttt{6b740dd} & Impayé antérieur saisissable à la migration. \\
\texttt{e1b2bd2} & Matricule de promotion ancré sur l'année et la date
  d'entrée. \\
\texttt{bd60d26} & Parcours scolaire, base des anciens élèves, sort des
  sortants. \\
\texttt{ec2c6a5} & Carte « Santé de la migration ». \\
\bottomrule
\end{tabularx}
\end{center}

\vfill
\noindent{\small\color{sagigray}Document généré le 27 juillet 2026 ---
SAGI SCHOOL by HADY GESMAN}

\end{document}
